\documentclass[10pt,twocolumn]{article}

\usepackage[letterpaper,margin=0.72in]{geometry}
\usepackage{booktabs}
\usepackage{array}
\usepackage{tabularx}
\usepackage{hyperref}
\usepackage{url}
\usepackage{xcolor}
\pagestyle{plain}

\title{Paper Agent: A Traceable Multi-Agent System for Verified Scholarly Paper Discovery}
\author{MON AI Team\\AI Agent-Based Domain Problem Solving Project\\Sungkyunkwan University}
\date{}

\begin{document}
\maketitle

\begin{abstract}
Academic paper discovery requires traceable retrieval, DOI integrity, metadata normalization, venue-policy enforcement, access-status inspection, and reproducible evaluation. We present \textit{Paper Agent}, an AI Agent-based literature-discovery and evaluation system implemented as a 12-stage multi-agent and tool-use pipeline. The workflow combines configured scholarly retrieval, Crossref verification, Unpaywall lookup, Cloudflare Worker execution, D1 traces, artifact storage, ranking, critic review, and dashboard reporting.

We evaluate the system through \textit{Benchmark v3}, a six-layer, 30-metric automated evaluation framework. Layers 1--4 and Layer 6 are computed from existing artifacts. Overall readiness is \textbf{PASS WITH CLAIM BOUNDARIES}. Baseline comparison is restricted to the T001--T003 common-support subset; T004--T020 remain artifact-level validation tasks, and T007 is marked \texttt{proposed\_agent\_missing}. Layer 5A is a quota-limited partial implementation audit with 22/125 successful rows (17.6\%); its evaluated subset contains no Proposed Agent rows, so the Proposed Agent Layer 5 score is \texttt{not\_available\_in\_subset}. Layer 5B provides deterministic semantic proxies for 125 rows, but supplements rather than replaces LLM or human semantic evaluation. The framework supports reproducible, claim-boundary-aware reporting; it does not support full T001--T020 comparative superiority or full semantic-quality validation claims.

\end{abstract}

\noindent\textbf{Keywords}---multi-agent systems, scholarly search, literature review automation, DOI verification, tool use, reproducibility.

\section{Introduction}
Literature review is a recurring bottleneck for business-school researchers. A defensible review must identify relevant work, verify metadata, retain DOI evidence, comply with an explicit journal policy, inspect access paths, and explain why papers were recommended or excluded. A fluent single-LLM answer can be useful, but it may not expose retrieval provenance, policy enforcement, DOI checks, or ranking reasons.

Paper Agent addresses this traceability gap through a multi-agent, retrieval-augmented, and tool-using workflow. Its stages expose records and failure boundaries rather than treating one model response as the final result. The project applies Multi-Agent decomposition, RAG-style grounding, MCP/Tool Use inspection surfaces, and benchmark evaluation.

Evaluation is part of the domain problem. A paper-discovery agent should be evaluated with preserved task scope, provenance, baseline availability, metadata checks, timeout evidence, semantic-audit coverage, and claim boundaries. Benchmark v3 separates controlled common-support comparison, artifact-level validation, quota-limited semantic audit, deterministic semantic proxies, and robustness indicators.

This project contributes: (1) a traceable 12-stage Paper Agent pipeline; (2) Benchmark v3, a six-layer and 30-metric automated framework; (3) reproducible validation artifacts; and (4) claim-boundary-aware reporting that prevents partial evidence from becoming a full-performance claim.

\section{Related Work}
\subsection{LLM-Based Scholarly Search}
LLM-based assistants can summarize papers, reformulate questions, and generate candidate reading lists. Their usefulness depends on grounding and metadata verification. Large scholarly corpora such as S2ORC show the scale and structure of the evidence that research-support systems must handle \cite{lo2020s2orc}. Ungrounded recommendations are risky in scholarly settings because a fabricated title, DOI, venue, or author list can consume researcher time and contaminate a review. Paper Agent therefore separates recommendation from verification and stores evidence for downstream inspection.

\subsection{Multi-Agent and Tool-Using Systems}
Multi-agent systems decompose complex tasks into role-specific components such as planners, executors, critics, and reporters. AutoGen demonstrates how multi-agent conversation can structure LLM application workflows \cite{wu2023autogen}. The value of decomposition is not the number of named agents alone; it is the ability to expose intermediate decisions and isolate failures. Paper Agent applies this principle to scholarly discovery by recording stage-level traces and preserving partial outcomes when a later service fails.

\subsection{Retrieval-Augmented Workflows}
Retrieval-augmented generation grounds model behavior in external evidence \cite{lewis2020rag}. Paper Agent follows the same grounding principle but emphasizes structured scholarly metadata and policy enforcement. Search results are filtered through an internal business-school journal allowlist and enriched through external metadata services before ranking.

\subsection{Benchmarking and Reproducibility}
Evaluation of agentic systems requires more than a single aggregate score. AgentBench provides an example of evaluating LLMs as agents across multiple environments and highlights the importance of explicit evaluation design \cite{liu2023agentbench}. Reproducible evaluation should disclose task scope, baselines, metrics, failure evidence, artifact locations, and incomplete runs. Paper Agent distinguishes controlled benchmark evidence from artifact-only expansion evidence.

\section{Method and System Design}
Paper Agent is a multi-agent decision-support pipeline rather than a single LLM wrapper. Each stage records an inspectable output. The current Worker implementation uses Web of Science as the primary configured scholarly search provider and OpenAlex as a fallback for testing and quota resilience. Crossref performs DOI and metadata verification, while Unpaywall records open-access status. Cloudflare Workers execute the pipeline and Cloudflare D1 persists jobs, papers, scores, and traces.

\begin{table*}[t]
\centering
\footnotesize
\begin{tabularx}{\textwidth}{c l X X X}
\toprule
Stage & Agent Role & Input & Tool / Data Source & Output Trace \\
\midrule
1 & Planner & question, constraints & request payload & normalized scope \\
2 & Journal Selector & scope, field & internal allowlist & journal pool \\
3 & Search / Retriever & keywords, years & WoS primary; OpenAlex fallback & candidates \\
4 & Journal Filtering & candidates & internal allowlist & accepted/excluded rows \\
5 & Verifier & DOI metadata & Crossref & verification reasons \\
6 & Open Access & verified DOI & Unpaywall & OA status and URLs \\
7 & Storage & records, artifacts & D1, R2, optional Drive path & persistence trace \\
8 & Evaluation Prep & enriched records & scoring preparation & normalized inputs \\
9 & Relevance & metadata text & metadata scoring; optional semantic path & relevance components \\
10 & Ranking & score components & deterministic ranking & ordered results \\
11 & Critic & ranked results & rule-based default; opt-in LLM path & critic flags \\
12 & Report / Delivery & results, traces & report generator, dashboard & downloadable outputs \\
\bottomrule
\end{tabularx}
\caption{Twelve-stage Paper Agent workflow. Each stage has an explicit trace boundary.}
\label{tab:pipeline}
\end{table*}

\subsection{Traceability and Failure Isolation}
The pipeline records job progress and stage-level traces in D1. This preserves observable partial state: retrieval failure, timeout, OA lookup failure, or report-generation limitations can be recorded without silently rewriting earlier evidence. The staged benchmark wrappers similarly write isolated artifact directories and stop later batches after a task-level failure.

\subsection{Journal Policy and Verification}
The internal business-school allowlist is a deliberate constraint. It helps the system focus on journals recognized by the project domain, but it is not a universal quality oracle. The Journal Selector uses the selected business field and tier priorities to construct a search and filtering policy. The filtering stage then records whether a candidate remains eligible. This separation avoids presenting a venue-policy decision as a relevance judgment.

For eligible candidates, the Verifier uses Crossref to inspect DOI and metadata consistency. The pipeline records whether title, year, and journal metadata match sufficiently for downstream review. A verification outcome does not prove substantive quality; it reduces bibliographic uncertainty. The Open Access stage separately uses Unpaywall to record whether an OA PDF or landing page is available. Access status is retained as metadata rather than treated as an academic-quality signal.

\subsection{Scoring, Ranking, and Critic Review}
Ranking combines observable components such as relevance, journal fit, verification status, citation metadata, recency, and access metadata. The implementation retains metadata-oriented scoring as the inspectable default. Semantic relevance and LLM-backed critic analysis remain optional paths because external model latency and resource constraints can affect live execution. The Critic stage does not overwrite the retrieved record. It adds review flags and summaries so that uncertainty remains visible.

This design supports a conservative interpretation of automation. A ranking score is a prioritization aid, not a scholarly verdict. Researchers can inspect the score components, metadata evidence, and trace history when deciding whether to read, cite, exclude, or re-query a paper.

\subsection{Storage, Reporting, and Dashboard Surfaces}
Cloudflare D1 stores jobs, papers, evaluations, and stage traces. Cloudflare R2 is used for report-output persistence where configured. The report pipeline exposes CSV, Markdown, XLSX, and PDF outputs, while the dashboard presents research, operations, and evaluation views. The dashboard also distinguishes live D1 evidence, controlled benchmark evidence, legacy partial artifacts, mock blueprint elements, simulation panels, and planned-only features. These labels prevent interface polish from being mistaken for implementation or evaluation evidence.

\subsection{Why the Architecture Is Agentic}
The system is not a thin interface around one prompt. It combines role decomposition, retrieval tools, explicit verification, persistent traces, deterministic scoring, a critic stage, and report delivery. LLM-backed behavior is optional where appropriate; the default pipeline remains inspectable even when resource-intensive model calls are disabled.

\section{Implementation Architecture}
\subsection{Request Lifecycle}
A Paper Agent request starts with a research keyword, optional year boundaries, a result-size limit, and an optional journal-field selection. The Worker creates a search job immediately and exposes a job identifier to the dashboard. This asynchronous boundary is important: the browser does not need to keep a single request open while retrieval and enrichment continue. The dashboard polls for observable progress and can present the latest completed state, while the Worker advances the job through its stage traces.

The Planner records normalized runtime constraints. The Journal Selector resolves the field context and internal allowlist priorities. The Retriever then invokes the configured scholarly backend. In the deployed implementation, Web of Science is the primary provider, and OpenAlex is retained as a fallback for testing and quota resilience. Provider fallback is recorded in traces so that downstream users can distinguish a primary-provider result from a fallback result.

After retrieval, journal filtering applies the approved business-school allowlist. This stage intentionally separates source retrieval count from accepted count. A low accepted count is not necessarily a search failure: it may indicate that the source returned papers outside the policy pool. The Verifier then calls Crossref for DOI-based metadata enrichment. Unpaywall lookup follows as a separate access-status stage. Separating these stages prevents an OA failure from being misreported as a DOI failure or a journal-policy failure.

\subsection{Trace and Persistence Model}
D1 stores search jobs, candidate papers, evaluations, and agent traces. A trace includes stage order, stage identifier, agent name, status, summary, detail, input count, output count, timestamps, and optional error messages. This representation serves two audiences. Researchers can inspect the human-readable summaries, while operators can use the counts and timestamps to diagnose a stall or external-service failure.

The trace model also supports honest incompleteness. A pipeline does not need to erase a partial result when a later stage fails. For example, a retrieval timeout can be represented as a task-level failure while other tasks in an isolated batch retain their artifacts. This design differs from a product dashboard that displays only a final success badge. Paper Agent treats the path to a result as part of the result.

The benchmark schema has been extended with fields intended for batch orchestration, including parent run identifiers, batch identifiers, derived-run markers, merge status, retry count, and task-level error fields. However, the current staged artifact wrappers do not persist these batch fields to Production D1. This limitation is explicitly retained in the manuscript because schema availability and end-to-end persistence are different implementation claims.

\subsection{Evidence and Artifact Layers}
The project distinguishes several evidence layers. \textit{Live D1 verified} data describes runtime jobs and traces obtained from deployed services. \textit{Controlled T001--T003} evidence describes the audited quantitative comparison layer. \textit{Artifact-only dry-run} data describes isolated CSV outputs created through gated wrappers. \textit{Legacy partial artifacts} preserve earlier incomplete attempts and failure records. Finally, \textit{mock blueprint} and \textit{scenario simulation} panels support interface design but are not live benchmark evidence.

This taxonomy is enforced in both documentation and dashboard labels. The distinction matters because an interactive interface can look complete before every evaluation path has been audited. By separating evidence sources, the system avoids turning mock values, partial artifacts, or legacy runtime records into final claims.

\subsection{Output Surfaces}
The report pipeline exposes multiple artifact formats for different uses. CSV and XLSX files support analysis and inspection. Markdown supports readable review summaries. PDF supports distribution, subject to the current report-engine language constraints. Dashboard panels expose ranked papers, paper detail, workflow progress, system diagnostics, and available output links. R2-backed persistence is available for configured report outputs, while optional Drive-related paths remain distinct from the core metadata workflow.

Each output format should be interpreted as a view over stored evidence, not as an independent source of truth. A PDF can summarize a run, but it cannot establish that an unaudited benchmark scope was completed. The report pipeline therefore inherits the same claim boundaries as the underlying jobs, traces, and artifacts.

\subsection{Operational Safety Gates}
Benchmark expansion uses explicit wrappers rather than broad package commands. Wrappers expose plan and preflight modes, isolate artifact paths by run ID, refuse protected directories, reject broad scopes, require explicit acknowledgement flags for execution, and preserve local outputs for read-only inspection. The T007 timeout demonstrates why this gate matters. After Batch 1 produced a failure, the staged process stopped before T013--T018 and T019--T020. A one-shot run could have obscured the point at which resource pressure became significant.

The gate design does not solve every operational issue. It does not implement D1 batch-aware persistence, automatic resumability, or an approved retry policy. Instead, it creates a conservative foundation for later work. The current system favors evidence preservation over aggressive completion.

\section{Benchmark Design}
Benchmark v3 evaluates inspectable literature-discovery evidence across T001--T020. It defines six layers with five metrics each, for a total of 30 metrics.
\begin{table*}[t]\centering\footnotesize
\begin{tabularx}{\textwidth}{l l X}\toprule
Layer & Metrics & Interpretation \\\midrule
1. Foundation \& Reproducibility & manifest, parity, claim boundary, trace, critic coverage & evidence reproducibility \\
2. Schema \& Metadata & normalization, completeness, DOI format, parsing, mismatch & structured artifact quality \\
3. Deterministic Validity & DOI match, existence, journal policy, journal precision, OA & validity signals \\
4. Retrieval Accuracy & Precision@5, NDCG@5, Recall@20, hit rate, MRR & retrieval and ranking \\
5. Semantic Quality & relevance, constructs, context, confidence, reasoning & Layer 5A audit and Layer 5B proxy \\
6. Robustness \& Risk & distractors, hallucination, timeout, latency, cost & operational risk \\
\bottomrule\end{tabularx}
\caption{Benchmark v3: six layers and 30 metrics.}\label{tab:v3layers}\end{table*}

\subsection{Baseline Support Matrix}
Only T001--T003 contain Paper Agent, Rule-Based, Single LLM, and gold-target support simultaneously. They form the controlled common-support subset. T004--T020 remain artifact-level validation tasks unless baseline parity is proven. T007 is \texttt{proposed\_agent\_missing}.

\subsection{Layer 5A and Layer 5B}
Layer 5A prepared 125 judge rows but evaluated 22 successfully (17.6\%) under quota limits. The successful subset contains no Proposed Agent rows, so it is a partial implementation audit rather than a representative semantic-quality estimate. Layer 5B computes deterministic semantic proxies for 125 rows: query-title overlap, query-notes overlap, construct coverage, context match, and evidence-field presence. These proxies supplement but do not replace LLM or human semantic evaluation.

\subsection{Claim Boundary}
Benchmark v3 is \textbf{PASS WITH CLAIM BOUNDARIES}. It supports a reproducible automated framework, not full T001--T020 comparative superiority or full semantic-quality validation.

\section{Evaluation Protocol and Claim Boundaries}
\subsection{Why the Benchmark Is Layered}
Benchmark v3 is organized as a diagnostic framework rather than a single leaderboard. A scholarly-discovery workflow can fail in several distinct ways. It can return an irrelevant paper, attach malformed metadata, provide a plausible but mismatched DOI, violate an explicit journal policy, omit evidence needed for an audit, time out under external-service pressure, or summarize partial data too confidently. A single aggregate score would hide these failure modes. The six-layer structure therefore preserves a separate interpretation boundary for reproducibility, metadata, deterministic validity, retrieval accuracy, semantic quality, and robustness.

Layer 1 checks whether evidence can be reconstructed from tracked manifests, scripts, traces, critic outputs, and claim-boundary records. These checks are prerequisites for analysis: a ranking score is difficult to defend if the input scope or generation path cannot be recovered. Layer 2 measures schema and metadata quality. Normalization and completeness remain separate because a syntactically normalized row may still omit fields needed for scholarly review. Layer 3 reports deterministic validity signals. DOI exact matching, existence evidence, journal-policy indicators, and access-status records help distinguish bibliographic confidence from topical relevance.

Layer 4 measures retrieval behavior through Precision@5, NDCG@5, Recall@20, hit rate, and mean reciprocal rank. These metrics expose different properties: top-result overlap, ranking order, broader target coverage, binary recovery, and earliest-hit position. Layer 5 handles semantic quality cautiously by separating a quota-limited LLM-judge audit from deterministic semantic proxies. Layer 6 reports robustness and risk. Hallucination, timeout, latency, and cost-related indicators matter because a literature-review assistant must remain usable under live service constraints.

\begin{table*}[t]
\centering
\footnotesize
\begin{tabularx}{\textwidth}{l X X}
\toprule
Evidence class & Permitted interpretation & Unsupported interpretation \\
\midrule
T001--T003 common support & Cross-method retrieval comparison within three controlled tasks & Generalization to every T001--T020 task \\
T004--T020 artifact-level validation & Artifact generation, schema inspection, deterministic checks, and visible failures & Baseline-parity comparison without matching baseline support \\
Layer 5A evaluated subset & Judge-path implementation behavior for 22 quota-limited rows & Representative semantic estimate or Proposed Agent semantic score \\
Layer 5B deterministic proxy & Reproducible supplementary semantic indicators for 125 rows & Replacement for human or LLM semantic evaluation \\
Layer 6 robustness evidence & Observed operational risk in retained artifacts & Proof that runtime limits have been resolved \\
\bottomrule
\end{tabularx}
\caption{Evidence classes and claim boundaries in Benchmark v3.}
\label{tab:evidenceboundary}
\end{table*}

\subsection{Baseline Support Protocol}
The Baseline Support Matrix is generated from existing artifacts. For each task, it records whether Proposed Agent outputs, Rule-Based outputs, Single LLM outputs, and gold or target evidence are available. A task is comparable only when the required support exists. This prevents a common evaluation error: treating the presence of one method's artifact as if it implied parity across every method. Under the current matrix, T001--T003 form the common-support subset. T004--T020 remain artifact-level validation tasks unless baseline parity is proven. T007 is explicitly marked \texttt{proposed\_agent\_missing}.

This distinction constrains retrieval interpretation. The controlled T001--T003 table is useful because the compared methods share a bounded evaluation context. It is not a universal leaderboard. The Single LLM baseline has higher Precision@5 and NDCG@5 in this subset, while Paper Agent retains explicit journal-policy enforcement, DOI verification, traceability, and visible failure reporting. The project does not collapse these different properties into an unsupported global ordering.

Gold or target labels are also bounded references rather than complete descriptions of the literature. A relevant paper can be absent from a finite target set. Retrieval scores therefore measure overlap with current target evidence. DOI, metadata, policy, and trace metrics provide complementary signals. Future work should expand DOI-backed targets with human audit while preserving the distinction between target coverage and universal relevance.

\subsection{Layer 5A Representativeness Audit}
The Layer 5A input population contains 125 prepared judge rows. Quota pressure allowed 22 rows to be evaluated successfully, which is 17.6\% coverage. The representativeness audit compares the full input population with the successful subset by task identifier, method, and candidate rank where available. Its purpose is not to repair missing coverage statistically. Instead, it discloses the degree to which the available subset differs from the intended evaluation population.

The successful subset contains no Proposed Agent rows. Consequently, the Proposed Agent Layer 5 score is recorded as \texttt{not\_available\_in\_subset}. Assigning a score would fabricate evidence. The correct interpretation is narrower: the partial run confirms that the judge-path implementation can generate and retain some evaluated rows under quota constraints. It does not establish representative semantic quality for Paper Agent, any baseline, or the full dataset.

This boundary matters because semantic judgments can be especially sensitive to sample selection. Candidate rank, research method, query language, abstract availability, and metadata density may all affect the difficulty of judgment. A future semantic protocol should predefine sampling, stratify across tasks and methods, retain failed calls, and compare automatic judgments with a human-reviewed subset.

\subsection{Layer 5B Deterministic Semantic Proxy}
Layer 5B computes deterministic semantic proxies for all 125 prepared rows. Available fields support signals such as query-title overlap, query-abstract-or-notes overlap where text exists, construct-coverage indicators, method-context match indicators, and evidence-field presence. Because these computations are deterministic, they can be rerun from tracked artifacts without spending judge quota or modifying raw outputs.

The proxy has a deliberately limited role. A high title-overlap score can reveal lexical alignment, and an evidence-field-presence score can reveal whether a candidate is sufficiently documented for later review. Neither signal establishes conceptual relevance. A paper can use different terminology for the same construct, and a paper with matching terms can still be methodologically unsuitable. Layer 5B is therefore a supplementary observability mechanism, not a substitute for LLM or human semantic evaluation.

\subsection{Artifact Preparation and Normalization}
Benchmark v3 supplements consume retained artifacts rather than rerunning scholarly search. This constraint improves reproducibility: generated matrices and summaries can be regenerated from a stable snapshot without introducing changes from live API responses. Normalization shapes metadata fields, DOI strings, journals, statuses, ranks, and evidence text into consistent rows for deterministic validation.

Normalization does not silently repair uncertain metadata. A malformed DOI remains a validity signal. A missing candidate artifact remains a support gap. A judge failure remains part of the coverage record. This policy treats missing evidence as information rather than imputing an optimistic value. It also permits engineering improvements to be evaluated later against the same preserved artifact lineage.

\subsection{Operational Risk Protocol}
Operational metrics are derived from retained jobs and artifacts. Timeout rate and latency per task are reported because external retrieval, enrichment, and polling behavior affect practical usability. Hallucination rate remains visible because metadata-oriented safeguards can reduce but not eliminate bibliographic uncertainty. These indicators characterize whether the workflow can be trusted under constrained resources.

The current framework records that D1 batch-aware persistence remains incomplete. Existing artifacts support deterministic supplements, but artifact generation should not be described as complete Production D1 batch orchestration. Schema availability, local artifact presence, and end-to-end persistence are different claims. The same distinction applies to dashboard summaries: a polished interface is a view over evidence, not an independent validation source.

\section{Experiments}
Benchmark v3 normalizes existing artifacts across T001--T020 without rerunning scholarly search. Layers 1--4 and Layer 6 are computed. Cross-method comparison is restricted to T001--T003 common support; T004--T020 remain artifact-level validation, and T007 lacks a Proposed Agent artifact.

Layer 5A preserves the quota-limited judge subset and failures. Layer 5B computes deterministic proxies for all 125 prepared rows. No semantic scores are fabricated. D1 batch-aware persistence remains incomplete.

\section{Results and Discussion}
\begin{table}[t]\centering\footnotesize
\begin{tabular}{lccc}\toprule Metric & Rule & Single LLM & Paper Agent \\\midrule
Precision@5 & 0.1333 & 0.6667 & 0.1333 \\
NDCG@5 & 0.3579 & 0.9949 & 0.3579 \\
DOI Accuracy & 1.0000 & 1.0000 & 1.0000 \\
Top Journal Prec. & 1.0000 & 0.9333 & 1.0000 \\
Hallucination Rate & 0.0 & risk noted & 0.0 \\
\bottomrule\end{tabular}\caption{Controlled T001--T003 common-support comparison.}\label{tab:controlled}\end{table}
The Single LLM baseline has higher Precision@5 and NDCG@5 in T001--T003. Paper Agent is not claimed to outperform baselines globally. Its contribution is an inspectable tradeoff: policy filtering, DOI verification, traces, artifacts, and failure evidence.

\begin{table}[t]\centering\footnotesize
\begin{tabular}{lr}\toprule Metric & Value \\\midrule
Schema normalization & 1.0000 \\
Metadata completeness & 0.9854 \\
DOI format validity & 0.9678 \\
DOI exact match & 0.6930 \\
Paper existence & 0.6930 \\
Top journal precision & 0.8129 \\
Paper Agent Precision@5 (T001--T003) & 0.1333 \\
Paper Agent NDCG@5 (T001--T003) & 0.3579 \\
Paper Agent Recall@20 (T001--T003) & 0.5000 \\
Hallucination rate & 0.3070 \\
Timeout rate & 0.1111 \\
Latency per task & 204.60s \\
\bottomrule\end{tabular}\caption{Scoped Benchmark v3 indicators.}\label{tab:v3summary}\end{table}
The baseline-support matrix confirms 3/20 comparable tasks. Layer 5A evaluated 22/125 rows and includes no successful Proposed Agent rows. Layer 5B generated proxies for 125 rows. These supplements improve observability but do not establish semantic-quality superiority. The defensible result is \textbf{PASS WITH CLAIM BOUNDARIES}.


\subsection{Metadata and Validity Interpretation}
Schema normalization reaches 1.0000, and metadata completeness reaches 0.9854. These values indicate that retained artifacts are generally structured enough for automated analysis. They do not imply that every recommendation is correct. DOI format validity is 0.9678, while DOI exact-match rate and paper-existence rate are both 0.6930. The difference between syntactic DOI format and stronger existence-oriented evidence illustrates why the benchmark separates formatting from verification. A DOI-like string can be well formed without providing sufficient evidence for a confident recommendation.

Top journal precision is 0.8129 in the broader Benchmark v3 evidence. This metric reflects compliance with the internal business-school allowlist, not a universal measure of research quality. A journal outside the project list is not necessarily weak, and a journal inside the list does not guarantee that an individual paper is relevant. Separating journal policy, DOI evidence, and retrieval metrics reduces the chance that one institutional preference is mistaken for a complete relevance judgment.

\subsection{Retrieval Interpretation}
For T001--T003, Paper Agent records Precision@5 of 0.1333, NDCG@5 of 0.3579, and Recall@20 of 0.5000. The Single LLM baseline records higher Precision@5 and NDCG@5 in the controlled table. This result should be stated directly. Paper Agent does not dominate pure relevance-ranking metrics in the current common-support subset. Its current contribution lies in an inspectable workflow that combines retrieval with policy checks, verification, persistent traces, and conservative reporting.

The tradeoff is useful for system design. An evaluation that reports only overlap metrics can reward plausible relevance ranking while omitting evidence-integrity requirements. An evaluation that reports only DOI validity can ignore whether users receive useful papers. Benchmark v3 retains both classes of evidence. Future optimization should improve retrieval quality without weakening DOI, policy, and trace safeguards. A larger baseline-parity set is also necessary before making broader comparative statements.

\subsection{Semantic-Quality Interpretation}
Layer 5A returned 22 successful rows from 125 prepared inputs. This 17.6\% coverage is too limited for representative semantic-quality conclusions. More importantly, the successful subset contains no Proposed Agent rows. Reporting a Proposed Agent semantic score would fabricate evidence. The benchmark instead reports \texttt{not\_available\_in\_subset} and uses the representativeness audit to preserve the limitation.

Layer 5B broadens deterministic visibility across all 125 prepared rows. Its proxies can identify sparse evidence fields, weak lexical alignment, missing construct terms, or limited context matches. These signals are useful for debugging and future audit prioritization. They are not a semantic verdict. Scholarly relevance can depend on theoretical framing, research design, population, setting, and contribution in ways that lexical proxies cannot capture fully.

\subsection{Robustness and Risk Interpretation}
The hallucination rate of 0.3070 and timeout rate of 0.1111 remain material risks. Latency per task is 204.60 seconds. These observations constrain deployment claims. A traceable system can still be slow for some interactions, and verification infrastructure can still leave uncertain recommendations. The value of the current framework is that these issues remain measurable and visible. Future versions should reduce the rates through improved retrieval, bounded retries, caching, and expanded verification rather than removing inconvenient evidence from reports.

\begin{table*}[t]
\centering
\footnotesize
\begin{tabularx}{\textwidth}{l X X}
\toprule
Observed condition & Current interpretation & Future mitigation \\
\midrule
Only T001--T003 have common baseline support & comparative scope remains partial & collect parity-preserving baseline artifacts for additional tasks \\
T007 is \texttt{proposed\_agent\_missing} & missing artifact remains visible & bounded recovery with preserved failure evidence \\
Layer 5A covers 22/125 rows & quota-limited implementation audit only & approved judge budget and stratified human audit \\
No Proposed Agent rows in successful Layer 5A subset & semantic score unavailable for Proposed Agent & balanced method sampling and retained failed-call records \\
Hallucination rate is 0.3070 & metadata risk remains material & strengthen retrieval, verification, and uncertainty labels \\
Timeout rate is 0.1111 & runtime resilience remains incomplete & bounded retries, caching, and provider-aware backoff \\
\bottomrule
\end{tabularx}
\caption{Current limitations, conservative interpretations, and future mitigations.}
\label{tab:mitigation}
\end{table*}

\subsection{Reporting Discipline}
The promotion gate reports \textbf{PASS WITH CLAIM BOUNDARIES}. This phrase is not cosmetic. It communicates that deterministic layers and robustness metrics have been computed while semantic-audit coverage and baseline parity remain incomplete. The dashboard, final report, and presentation should retain the same qualifier. Generated summaries must not turn artifact availability into parity, proxy metrics into semantic ground truth, or partial judge coverage into representative evaluation.

This reporting discipline is a practical contribution of the project. Agent systems often combine live APIs, asynchronous jobs, generated artifacts, and optional model calls. Their evidence can be uneven. A claim-boundary-aware report gives developers and evaluators a structured way to communicate what is available, what failed, and what still needs review.

\section{Limitations and Ethics}
Baseline parity is partial. Layer 5A is quota-limited to 22/125 successful rows and contains no successful Proposed Agent rows. Layer 5B does not replace semantic evaluation. Full semantic-quality validation and full T001--T020 comparative superiority are not supported. The hallucination rate of 0.3070 and timeout rate of 0.1111 remain material risks. Broader evaluation requires funded judge coverage or human review.

The journal allowlist may create algorithmic gatekeeping. Ranking can amplify citation and hierarchy biases. DOI verification reduces bibliographic uncertainty but does not establish substantive quality. Human scholarly review remains the final authority.

\section{Reproducibility and Design Implications}
\subsection{Reproducible Artifact Chain}
The repository is designed so that an evaluator can inspect the path from task metadata to normalized rows, metric outputs, supplement reports, promotion-gate summaries, and presentation language. The Baseline Support Matrix records where comparison is possible. The Layer 5A representativeness report records the gap between prepared judge inputs and successful evaluations. The Layer 5B proxy report records deterministic semantic signals. The robustness layer records timeout, hallucination, and latency observations. Together, these artifacts make uncertainty inspectable.

This artifact chain is important for live-service systems. Search providers and metadata APIs can change over time, and a rerun can produce different candidates or response timing. A reproducible supplement should therefore preserve the artifact snapshot used for analysis and derive reports from that snapshot. Reproducibility does not mean that every future live call is identical. It means that the evidence supporting a submitted claim can be reconstructed and audited.

\subsection{Claim Boundaries as a System Requirement}
Claim boundaries are not a final editing pass. They influence schema design, metric generation, dashboard labels, and presentation structure. A task-level missing artifact must remain visible in the support matrix. A quota-limited judge run must expose coverage. A deterministic proxy must retain its supplementary label. A promotion gate must communicate its qualifier. These requirements reduce the risk that downstream summaries accidentally overstate incomplete evidence.

The approach is especially relevant for agentic workflows. Multi-stage systems can appear more capable than their least reliable dependency. A retrieval provider may return candidates while a metadata service fails. An artifact directory may exist while batch-aware D1 persistence remains incomplete. A dashboard may render a summary while baseline parity is unavailable. Treating these states separately is necessary for honest engineering and scholarly use.

\subsection{Implications for Literature-Review Assistance}
Paper Agent illustrates a broader design principle: literature-review automation should be optimized for evidence quality as well as ranking quality. Relevance remains essential, and the current T001--T003 comparison shows that Paper Agent needs retrieval improvement. At the same time, a useful research assistant should expose metadata integrity, policy decisions, access status, trace history, and failure reasons. Researchers need enough context to challenge a recommendation rather than merely accept it.

A practical deployment should support iterative review. Users may inspect a recommended paper, reject a venue-policy decision, refine a query, or request broader coverage. Future feedback mechanisms should preserve provenance and avoid automatically promoting unreviewed feedback into gold labels. Human-audited governance is particularly important when the system is used in academic settings where citation choices influence literature visibility.

\subsection{Future Evaluation Protocol}
The next evaluation cycle should expand common-support baseline artifacts beyond T001--T003, recover T007 with a bounded retry policy, and obtain balanced semantic judgments across tasks and methods. A paid judge budget can reduce quota pressure, but the protocol should still use stratified sampling and retain failed-call evidence. Human reviewers should examine a subset of semantic judgments to estimate agreement and identify construct-level errors that deterministic proxies cannot detect.

Future work should also connect artifact batches to Production D1 with explicit parent-run, batch, retry, and error fields. That implementation would improve trace continuity from staged execution to dashboard reporting. It should be introduced with migration freeze rules, controlled backfill, and regression checks for existing evidence. Until then, the manuscript retains the narrower distinction between artifact-level validation and end-to-end batch-aware persistence.

\section{Conclusion}
Paper Agent combines a traceable multi-agent pipeline with Benchmark v3, a six-layer and 30-metric automated framework. The defensible conclusion is \textbf{PASS WITH CLAIM BOUNDARIES}. Layers 1--4 and Layer 6 are computed; comparison is restricted to T001--T003 common support; Layer 5A is a quota-limited partial implementation audit; and Layer 5B is supplementary. The benchmark does not support full T001--T020 comparative superiority or full semantic-quality validation.

\section{Reproducibility Statement}
The repository contains code, prompts, task metadata, gold targets, normalization scripts, baseline-support outputs, Layer 5A representativeness artifacts, Layer 5B proxy artifacts, robustness summaries, and claim-boundary documentation. Supplements are derived from existing artifacts without rerunning scholarly search or the judge. Raw judge outputs remain preserved.

\section{Future Work}
Future work should expand baseline parity beyond T001--T003, recover T007 under a bounded retry policy, obtain broader semantic coverage through a funded judge budget or human evaluation, connect batch-aware D1 persistence, strengthen timeout handling, and expand DOI-backed gold targets.

\section{Implementation and Evaluation Checklist}
Before final submission, the team should verify generated supplement artifacts, retain the T001--T003 common-support boundary, and disclose Layer 5A quota limits. A dashboard or report must not be used as evidence of full comparative or semantic validation unless the underlying scope supports that claim.


The final repository audit should confirm that baseline-support outputs, Layer 5A representativeness outputs, Layer 5B proxy outputs, and promotion-gate summaries are versioned together. Readers should be able to reconstruct the evidence chain from task metadata to normalized rows, metric summaries, and reporting language. This requirement is stronger than a screenshot-based demo check because reproducibility depends on tracked scripts and artifacts.

\newpage

\begin{thebibliography}{4}
\bibitem{lo2020s2orc} K. Lo et al., ``S2ORC: The Semantic Scholar Open Research Corpus,'' in \textit{Proceedings of ACL}, 2020.
\bibitem{wu2023autogen} Q. Wu et al., ``AutoGen: Enabling Next-Gen LLM Applications via Multi-Agent Conversation,'' arXiv:2308.08155, 2023.
\bibitem{lewis2020rag} P. Lewis et al., ``Retrieval-Augmented Generation for Knowledge-Intensive NLP Tasks,'' in \textit{Advances in Neural Information Processing Systems}, 2020.
\bibitem{liu2023agentbench} X. Liu et al., ``AgentBench: Evaluating LLMs as Agents,'' arXiv:2308.03688, 2023.
\end{thebibliography}

\end{document}
